% 03_models.tex - the Raven model and the existing MODFLOW 6 model
\section{The two models}

\step{Copy the start folder and look at the MODFLOW model}
\begin{shell}
$ cp -r $WS/model/start $WS/work
$ cd $WS/work
\end{shell}
The folder holds the Raven model of Tutorial Creek without groundwater, the HRU polygons \cmd{hrus.geojson}, and the
MODFLOW~6 model in the folder \cmd{model}:
\term{t4_ls_model.txt}
The name file lists the packages of the model:
\term{t4_nam.txt}
\why the model has:
\begin{itemize}
\item A regular grid (package \cmd{DIS}) of \V{t4:nrow}~rows and \V{t4:ncol}~columns of \V{t4:cellft}-ft cells
  (\V{t4:cellm}~m) and two layers, with lengths in feet. Its origin is given in feet, in UTM zone 11N:
\term{t4_dis_head.txt}
\item A steady-state period of one day, then monthly periods, \V{t4:nper} in all, from 1 January 2010:
\term{t4_tdis_head.txt}
\item Its own recharge (\cmd{RCH\_MODEL}) and ET (\cmd{EVT\_MODEL}), given as arrays.
\item \V{t4:nriv} river cells (\cmd{RIV\_CREEK}). Their auxiliary variable \cmd{SUBBASIN} (an extra column) names the Raven
  subbasin of each river reach:
\term{t4_riv_head.txt}
\item A pumping well (\cmd{WEL\_PUMP}, monthly rates) and a head boundary at the outlet (\cmd{GHB\_OUTLET}, \V{t4:nghb}
  cells).
\end{itemize}
